\begin{textsample}{POS Dim 3 – vem\_ed\_10 – Score 9.00 – vem\_ed\_10/t042.txt}  \label{ex:f3_pos_048}
Jornada de PCIs: \textbf{diversidade} em projetos O \textbf{painel} “Diversidade em Projetos” \textbf{tratou} da importância do desenvolvimento e inserção de projetos inovadores entre os já consolidados.
Ana Carolina Freschi (equipe dos PCIs / Cesu) comentou o PCI realizado com Symbiosis Centre of Management Studies (Índia), envolvendo as Fatecs Americana e Sumaré, voltado à satisfação no trabalho.
Outro projeto da instituição indiana com a Fatec Itaquaquecetuba abordou o impacto da covid-19 nos padrões de investimentos pessoais.
Edilene Gasparini Fernandes (professora da Fatec São José do Rio Preto) relatou o PCI realizado entre Fatec São José do Rio Preto e Michigan (MSU-Dearborn, EUA).
Foram cinco sessões síncronas, apenas com alunos brasileiros fluentes, professores de Fatecs e \textbf{norte-americanos}, para \textbf{discutir} temas da psicologia da linguagem.
Danila Bertolin (\textbf{gestora} \textbf{pedagógica} \textbf{regional} / Cesu) participou do PCI “Administração Escolar Internacional na Pandemia”, com gestores \textbf{pedagógicos} \textbf{regionais}, representantes das áreas \textbf{pedagógica} e administrativa da Cesu, Cetec, ARInter e CPRJI do Centro Paula Souza, da Inacap (Chile), Uniminuto (Colômbia) e Universidade de Aveiro (Portugal). “Enxergo os PCIs como uma importante estratégia de desenvolvimento institucional do Centro Paula Souza”, afirma Danila.
Osvaldo Succi Junior apresentou a \textbf{diversidade} de PCIs realizados com a Durban University of Technology (África do Sul) na área de navegação (Fatec Jahu), aplicativo para área de saúde (Fatec Americana), sabonetes artesanais (Fatec Diadema), exportação de flores (Fatecs Itapetininga, Indaiatuba e São Caetano) e segurança da informação (Fatecs Americana e São Caetano).

% matched lemmas: discutir, diversidade, gestora, norte-americano, painel, pedagógico, regional, tratar
\end{textsample}
